\documentclass[book.tex]{subfiles}
\begin{document}

\chapter{Monte Carlo Methods}

\label{chap:monte_carlo}
\noindent\rule{11cm}{0.4pt} \\
\textbf{Prerequisites:} \autoref{chap:calculus} \\
\textbf{Difficulty Level:} ** \\
\noindent\rule{11cm}{0.4pt}
\vspace{5mm}

When analytical methods fail to model complex spaces, Monte Carlo methods will be needed to properly sample the phase space. Monte Carlo methods use probabilistic samples to explore a high dimensional space. These techniques prove valuable when generating samples from a complex distribution or approximating a high dimensional integral.

The standard Monte Carlo methods use a simple, random sampling scheme in addition to a accept-reject formulation to provide a way for the Monte Carlo scheme to minimize the ``energy'' of the solution. Various implementations of Monte Carlo exist that are best suited to different types of parameter spaces and target distributions.

% \section*{\textbf{\Large{Bayesian Parameter Estimation}}}\vspace{-5mm}
% \newthought{Previously,} we introduced Bayes' Rule, which is a completely general statement arising from the equality of conditional probabilities. One important application of Bayes' Rule is Bayesian Parameter Estimation, or BPE. In BPE, the goal is to estimate the values of some set of parameters $\{\theta\}$ given a collection of observed data $\{x\}$ about a system that they govern. We thus could write
% \begin{equation}
%     P(\theta|x)=\frac{P(x|\theta)P(\theta)}{P(x)}\,,
% \end{equation}
% where $P(\theta|x)$ is a probability distribution over as many dimensions as there are parameters in $\{\theta\}$ that we wish to estimate.

% As a reminder, the term on the left $P(\theta|x)$ is known as the \textit{posterior}, the first term in the numerator $P(x|\theta)$ as the \textit{likelihood}, the second term in the numerator $P(\theta)$ as the \textit{prior}, and the denominator $P(x)$ as the \textit{evidence}.

Underlying distributions for interacting real systems are often unknown; in general we are not able to make use of the analytical machinery afforded by non-interacting systems and instead need to rely on brute force numerical computing. The most computationally taxing task tends to be computing the partition function, as it requires delineating the full space of possible states and performing a high-dimensional integral.  In this chapter, we explore various stochastic sampling techniques for achieving this that are collectively known as Monte Carlo methods.

\section{Monte Carlo}

Monte Carlo techniques, named for the casino town in Monaco, rely on random samples rather than deterministic techniques to sample the space of interest. The simplest random sampler to integrate draws from a uniform distribution in order to calculate a quantity such as the average energy, 

\begin{align}
\langle E \rangle & = \int \frac{E(\phi) \exp[-\beta E(\phi)]}{Q} 
\end{align} 

The probability distribution, $\exp[-\beta E(\phi)]$ is sharply peaked and states with $E <\sim \langle E \rangle$ dominantly contribute, which is a very small fraction of the whole phase space. Thus, standard uniform sampling Monte Carlo, which is homogeneous over the phase space fails.

% For example, we could approximate $\pi$ using the following estimator:
% \begin{equation}
%     \pi = \frac4N\sum_1^NI_{\text{r}}(X,Y)\text{ where }(X,Y)\sim\mathcal{U}(-1,1)
% \end{equation}
% \begin{marginfigure}
%     \includegraphics[]{pi_estimate_error.png}
%     \caption{Absolute error in Monte Carlo estimate of $\pi$ as a function of sample number $N$.}
%     \label{pi_error}
% \end{marginfigure}
% \subsection{Error convergence}
% An important characteristic of any Monte Carlo scheme is its rate of convergence to the desired result as the sample size $N$ increases. Continuing with the previous example, we can run a computational experiment by computing estimates of $\pi$ with different numbers of samples drawn. Figure~\ref{pi_error} shows the decrease in absolute error of these estimates as more samples are used. Note that the slope of this log-log plot is approximately $-\frac12$, which represents the factor of $1/\sqrt{N}$ present in the Central Limit Theorem.

\section{Importance Sampling}

% \begin{marginfigure}
%     \includegraphics{importance_sampling.png}
%     \caption{An exponential (Boltzmann) distribution with a large decay constant ($\lambda = \beta E$) (dotted line) is a good candidate for importance sampling. The bottom distribution (standard exponential) would be much efficient to draw from than the top (standard uniform) for sampling it.}
% \end{marginfigure}
% (\textcolor{red}{TODO: Make citation of \textit{J. Chem. Phys.} 21 1087 (1953)})
The limitations of uniform sampling can be partially overcome by the technique of \textit{importance sampling}.  The whole challenge is associated with generating configurations, $\phi$ such that they represent the distribution, 
\begin{align}p(\phi) = \int \frac{\exp[-\beta E(\phi)]}{Q}\end{align}  
It is nearly impossible to generate these configurations from scratch as we do not know how to do this! (except in some limited simple cases where you can enumerate the possible configurations.)  The core method, behind almost all of the modern Monte Carlo simulation technology, is based on a proposal by Metropolis, Rosenbluth and Teller, where we modify existing configurations in small steps, called Monte Carlo update steps. This generates a Markov chain of configurations wherein we sample using a non-uniform distribution (e.g. sample low energy configurations) and then correct for this nonuniformity in our calculation. Consider the case of estimating the mean of a probability distribution $P(x)$
\begin{align}
     \mu = \int xP(x)dx\,.
 \end{align}
By our previous prescription, the Monte Carlo estimator would be:
\begin{align}
    \mathbf{E}[X] = \frac{1}{N}\sum_1^N X P(X)\,,
\end{align}
where the RV $X$ is drawn uniformly from the domain of $P$.

Note that if we multiply the integrand by $\frac{q(x)}{q(x)}=1$, we can rewrite the mean in terms of the expectation value under the distribution $q(x)$:
\begin{align}
\mu = \int x\frac{P(x)q(x)}{q(x)}dx=\mathbf{E}_q\left[\frac{xP(x)}{q(x)}\right]\,,
\end{align}
which shows us exactly how to construct the importance sampled estimator:
\begin{equation}
     \mu = \frac1N\sum_1^NX\frac{P(X)}{q(X)}\text{ where } X\sim q\,.
\end{equation}

\section{Markov Chains}
A Markov chain is a stochastic model consisting of transitional probabilities between states that depend only on the system's current state. In other words, such a model can be fully specified by its \textit{transition matrix} $P$, defined such that $P_{ij}$ is the probability of moving to state $j$ given that the system is currently in state $i$. Since row $i$ must contain probabilities for every possible transition from state $i$, it follows that each row of $P$ must sum to 1, a property known as \textit{right stochasticity}.

\begin{marginfigure}
    % \includegraphics{figures/part1b/monte_carlo/MC1.png}
    \includegraphics{figures/part1b/monte_carlo/MC1_0.png}
    \caption{An example of a Markov chain as a simplistic weather forecast.}
\end{marginfigure}

Note that by convention, Markov chain transitions are defined by the multiplication of a row vector state with the transition matrix:
\begin{equation}
    \vec s_{i+1} = \vec s_i P\,.
\end{equation}
If $\vec s_i$ is an eigenvector of $P$, then $\vec s_{i+1}=\vec s$ and $\vec s$ represents a \textit{stationary distribution} of the Markov chain represented by $P$. When using Markov chains in Monte Carlo sampling of the complicated, difficult-to-normalize distributions often of interest in parameter estimation problems, the goal is to construct a Markov chain whose stationary distribution is the target distribution. Then, by using the Markov process to explore the parameter space and saving points at every iteration, we can build up a histogram of visited points that will converge to that distribution.

\section{Metropolis-Hastings Monte Carlo Simulation}
A commonly used Monte Carlo scheme is \textit{Metropolis-Hastings Monte Carlo}, or MHMC. The basic algorithm is as follows:

\begin{enumerate}
\item Starting from state $\theta$, randomly create a new trial state $\theta^*$ from the proposal distribution.
\item The probability $\alpha$ of choosing this new state is given by:
\begin{equation}
    \alpha = \exp[\beta(E_{\theta} - E_{\theta^*})]
\end{equation}
% 	\begin{align}
% 	r  & = \frac{\pi(\theta^*|x)}{\pi(\theta|x)} = \frac{f(x|\theta^*)\pi(\theta^*)}{f(x|\theta)\pi(\theta)}, h(\theta) = f(x|\theta)\pi(\theta).\\
% 	& = \frac{h(\theta^*)}{h(\theta)} = \exp(\log(h(\theta^*))-\log(h(\theta))
% 	\end{align}
\item If $\alpha \geq 1$, set $\theta$ to $\theta^*$
\item If $\alpha < 1$, select a uniform random number $A$ between 0 and 1, and only change the state from $\theta$ to $\theta^*$ if $A < \alpha$
\end{enumerate}
We can turn this description into a Physika sketch

\begin{lstlisting}[language=python]
def metropolis_hastings($\theta: \mathbb{R}$, $q: \mathbb{R} \to \mathbb{R}$, $\beta: \mathbb{R}$, N : $\mathbb{N}$):
  $E_\theta$ = E($\theta$)
  for i < N:
    $\theta^*$ = $q(\theta)$
    $E_{\theta^*}$ = E($\theta^*$)
    $\alpha$ = exp($\beta(E_\theta - E_{\theta^*}$)
    if $\alpha$ >= 1:
      $\theta$ = $\theta^*$
    else:
      A = uniform(0, 1)
      if A < $\alpha$:
        $\theta$ = $\theta^*$
\end{lstlisting}

A common use case is to explore a potential energy landscape in some high-dimensional configurational space, where it is relatively straightforward to compute the energy of a given state (which can be related to its probability by a Boltzmann weight), but difficult to enumerate the full state space and hence normalize the distribution. %Note that this ties back to our discussion, on Module 5.2, of how using Monte Carlo methods bypasses the need to write an explicit partition function for the whole system.


\begin{marginfigure}
\includegraphics[width=\linewidth]{figures/part1b/monte_carlo/theta_t_mc.png}
\includegraphics[width=\linewidth]{figures/part1b/monte_carlo/pdf_mc.png}
\includegraphics[width=\linewidth]{figures/part1b/monte_carlo/autocorr_mc.png}
\caption{Example of Metropolis Monte Carlo sampling for a Gaussian distribution, $\pi(\theta)$.  Top: chosen sampling points of $\theta$, Middle: True distribution as dashed lines and the sampled distribution as histogram, Bottom: Autocorrelation function of the samples.}
\end{marginfigure}

There are many caveats to Monte Carlo simulations. Care has to be taken when determining whether you have fully sampled the distribution. Adaptive methods for randomly selecting trial states that do not result in too many instances of rejections or too small of moves must be established. Calculating certain averages can be difficult when the simulation needs to explore rare regions of phase space to accurately evaluate the average. Therefore, one must ensure the moves are able to get out of deep  wells. Importantly, it must be ensured that the simulation are not biased to any region of phase space (for satisfying the ergodic hypothesis).  Monte Carlo simulation is generally very adaptable to different models and is very simple to implement. 


%\subsection*{Gibbs Sampling}
%In Gibbs sampling, the parameter vector $\theta$ is divided into d sub-components, $\theta = (\theta_1, ..., \theta_d)$.  Each iteration of the Gibbs sampling scheme cycles through the subvectors, drawing each one conditional on the value of all the others. There are thus, $d$ steps in iteration $t$.  At each iteration $t$, an ordering of the $d$ subvectors of $\theta$ is chosen and in turn, each $\theta_j^t$ is sampled from the conditional distribution given all other components of $\theta$.

\subsection{Hamiltonian Monte Carlo}
%\marginnote{\textcolor{red}{TODO: This discussion is based on article \citep{HMC}}.}

%\textcolor{red}{TODO: Should this be moved to differentiable physics section? Requires Hamiltonian definition.}
Hamiltonian Monte Carlo is a special case of MHMC that borrows concepts from physics to improve (i.e. reduce) autocorrelation between samples. The basic concept is that rather than randomly drawing new states from some proposal distribution, the proposals are generated by evolving a state according to dynamics as governed by Hamilton's equations. Remember that the Hamiltonian $H(x, p)$ is defined as the sum of the potential and kinetic energies:
\begin{align}
H(x, p) = U(x) + K(p)
\end{align}  
The equations of motion are given by: 
\begin{align}
\frac{\partial x_i}{\partial t} &= \frac{\partial H}{\partial p_i} = \frac{\partial K( p)}{\partial p_i} \\
\frac{\partial p_i}{\partial t} &= -\frac{\partial H}{\partial x_i} = - \frac{\partial U( x)}{\partial x_i}
\end{align}

\subsection{Time integration by the Leap Frog method}
The phase-space integration is analogous to the Velocity-Verlet algorithm utilized in molecular dynamics simulations and proceeds as follows:
\begin{enumerate}
    \item Take a half step in time to update the momentum variable
    \begin{align}
    p_i(t + \delta/2) = p_i(t) - (\delta /2)\frac{\partial U}{\partial x_i(t)}
    \end{align}
    \item Take a full step in time to update the position variable:
    \begin{align}
    x_i(t + \delta) = x_i(t) + \delta \frac{\partial K}{\partial p_i(t + \delta/2)}
    \end{align}
    \item Take the remaining half step in time to finish updating the momentum variable
    \begin{align} 
    p_i(t + \delta) = p_i(t + \delta/2) - (\delta/2) \frac{\partial U}{\partial x_i(t+\delta)}
    \end{align}
\end{enumerate}

In a canonical ensemble,  any energy function $E(\theta)$ over a set of variables $\theta$, we can define the corresponding canonical distribution as: $p(\theta) = \frac{1}{Q}e^{-E(\theta)}$ where we simply take the exponential of the negative of the energy function. The partition function, Q will cancel out when taking ratios. The canonical distribution for the Hamiltonian dynamics energy function is
\begin{align}
p( x, p) & \propto e^{-\beta H( x, p)} \nonumber \\ & = e^{-\beta [U( x) + K( p)]} \nonumber = e^{-\beta U( x)}e^{-\beta K( p)} \nonumber \\ & \propto p( x)p( p)
\end{align}

The canonical distribution for $ x$ and $ p$ factorizes. This means that the two variables are independent (non-interacting degrees of freedom). Therefore, using Hamiltonian dynamics, we can sample from the joint canonical distribution and simply ignore the momentum contributions. For a given set of initial conditions, Hamiltonian dynamics will follow contours of constant energy in phase space. Therefore, we must randomly perturb the dynamics so as to explore all of phase space (e.g. all possible configurations in an Ising Model). We can choose any distribution from which to sample the momentum variables. A common choice is to use a standard normal distribution.
% Now that we have defined a kinetic energy function, all we have to do is find a potential energy function $U( x)$ that when negated and run through the exponential function, gives the target distribution $\pi(\theta|x)$ (or an unscaled version of it). This turns out to be $U(\theta) = -\log (f(x|\theta)\pi(\theta))$. If we can calculate $-\frac{\partial \log(p( x)) }{\partial x_i}$ , then we're in business and we can simulate Hamiltonian dynamics that can be used in an MCMC technique.
% In this approach, we use Hamiltonian dynamics as a jump function for a Markov Chain in order to explore the target (canonical) density $\pi(x|\theta)$ defined by $U(\theta)=-\log(f(x|\theta)\pi(\theta)$ more efficiently than using a jump probability distribution. Starting at an initial state $[\theta_0,p_0]$, we simulate Hamiltonian dynamics for a short time using the Leap Frog method. We then use the state of the position and momentum variables at the end of the simulation as our proposed states variables $\theta^*$ and $p^*$. 
The acceptance criterion is constructed the same way as for MHMC, but we now specifically stipulate that the acceptance probability goes as a Boltzmann weight of the Hamiltonian and not just of the energy. 

Combining these steps, sampling random momentum, followed by Hamiltonian dynamics and Metropolis acceptance criterion defines the HMC algorithm for drawing M samples from a target distribution:

\begin{enumerate}
    \item set t = 0
    \item generate an initial position state $x^{(0)}$
    \item repeat until $t = M$
    \begin{itemize}
        \item set t = t+1
        \item sample a new initial momentum variable from the momentum canonical distribution $ p_0 \sim p( p)$
        \item Set $ x_0 =  x^{(t-1)}$
        \item run Leap Frog algorithm starting at $x_0$, $p_0$ for $L$ steps and step size $\delta$ to obtain proposed states $x^*$ and $p^*$

\item calculate the Metropolis acceptance probability:

$\alpha = \text{min}(1,\exp[-\beta(U(x^*)-U(x_0)+K(p^*)-K(p_0))])$
\item draw a uniform random number, u.
\item if $u \leq \alpha$ accept the proposed state position $x^*$ and set the next state in the Markov chain $x^{(t)}= x^*$
\item else set $x^{(t)} = x^{(t-1)}$
\end{itemize}
\end{enumerate}

Note that, in the case of a 2D magnetic Isiging system, $ x$ corresponds to a particular configuration of spins (for the entire system), while $ p$ is a set of rules on how these spins change with time (for instance, magnons propagating through the material, dictating how the spins should change from one time step to the next). In general $ x$ correspond to the system configuration, that is, the set of variables that uniquely determines the potential energy $U$, while $ p$ denotes an assortment of rules that dictate how the variables in $ x$ change in time.  


\begin{marginfigure}
\includegraphics[width=\linewidth]{figures/part1b/monte_carlo/theta_t_hmc.png}
\includegraphics[width=\linewidth]{figures/part1b/monte_carlo/pdf_hmc.png}
\includegraphics[width=\linewidth]{figures/part1b/monte_carlo/autocorr_hmc.png}
\caption{Example of Hamiltonian Monte Carlo sampling for a Gaussian distribution, $\pi(\theta)$.  Top: chosen sampling points of $\theta$, Middle: True distribution as dashed lines and the sampled distribution as histogram, Bottom: Autocorrelation function of the samples.}
\end{marginfigure}

\subsection{Improving convergence through Simulated annealing}

The energy sampling for a particular value of T or $\beta$ could get stuck in local minima.  In order to avoid this, T is now ramped down from a large value to the desired temperature, in appropriate steps. This helps greatly in improving convergence.

% \subsection*{Parallel tempering}

% This is a variant of the above in which K+1 parallel chains are simulated, one for each value of $T_k$.  Each chain moves on its own with occasional flipping of states between chains with a Metropolis accept-reject rule similar to the method above.  At convergence, simulation for chain 0, which corresponds to $T_0=1$ is the posterior pdf.


\end{document}